\par\Needspace{8\baselineskip}\bigskip
\pdfbookmark[3]{Exposition, 11}{pallas-II-article11}
\begin{center}11.\end{center}
\noindent\hspace*{\parindent}Les équations 42, 43 nous donnent celle-ci:
\[
\dfrac{x}{r}\cos\AscendingNode+\dfrac{y}{r}\sin\AscendingNode=\cos(v-\AscendingNode).
\]
En la différentiant nous trouvons
\[
\begin{aligned}
&\left(\mathrm{d}x-\dfrac{x\mathrm{d}r}{r}\right)\dfrac{\cos\AscendingNode}{r}+\left(\mathrm{d}y-\dfrac{y\mathrm{d}r}{r}\right)\dfrac{\sin\AscendingNode}{r}\\
&\quad-\left(\dfrac{x}{r}\sin\AscendingNode-\dfrac{y}{r}\cos\AscendingNode+\sin(v-\AscendingNode)\right)\mathrm{d}\AscendingNode\\
&\quad+\sin(v-\AscendingNode)\mathbin{.}\mathrm{d}v=0
\end{aligned}
\]
ou après les substitutions convenables
\[
\begin{aligned}
&(Bz-Cy)\dfrac{\cos\AscendingNode}{r^3}\mathrm{d}t+(Cx-Az)\dfrac{\sin\AscendingNode}{r^3}\mathrm{d}t\\
&\quad-(1-\cos i)\sin(v-\AscendingNode)\mathbin{.}\mathrm{d}\AscendingNode+\sin(v-\AscendingNode)\mathbin{.}\mathrm{d}v=0.
\end{aligned}
\]
\par Nous avons donc, en tirant les valeurs de \(Bx-Cy\), \(Cx-Az\) des équations 48, 49,
\[
\begin{array}{rl}
67)\quad&\mathrm{d}v=\dfrac{\BareRoot{mp}}{rr}\mathrm{d}t+(1-\cos i)\mathrm{d}\AscendingNode.
\end{array}
\]
Différentiant l’équation \(\dfrac{p}{r}=1+e\cos(v-\tilde\omega)\), on trouve
\[
\begin{aligned}
\dfrac{\mathrm{d}p}{r}-\dfrac{p\mathrm{d}r}{rr}&=\cos\varphi\mathbin{.}\cos(v-\tilde\omega)\mathbin{.}\mathrm{d}\varphi\\
&\quad-e\sin(v-\tilde\omega)\mathbin{.}\mathrm{d}v+e\sin(v-\tilde\omega)\mathbin{.}\mathrm{d}\tilde\omega.
\end{aligned}
\]
\par Substituant d’abord pour \(\mathrm{d}r\), \(\mathrm{d}v\) leurs valeurs tirées des équations 58, 67, on a
\[
\begin{aligned}
\dfrac{\mathrm{d}p}{r}&=\cos\varphi\mathbin{.}\cos(v-\tilde\omega)\mathbin{.}\mathrm{d}\varphi\\
&\quad+e\sin(v-\tilde\omega)\mathbin{.}(\mathrm{d}\tilde\omega-(1-\cos i)\mathrm{d}\AscendingNode)
\end{aligned}
\]
et ensuite, mettant pour \(\mathrm{d}p\), \(\mathrm{d}\varphi\) leurs valeurs:
\[
\begin{aligned}
&e\sin(v-\tilde\omega)\mathbin{.}(\mathrm{d}\tilde\omega-(1-\cos i)\mathrm{d}\AscendingNode)\\
&\quad=-aa\cos\varphi\mathbin{.}\cos(v-\tilde\omega)\mathbin{.}\sin(v-\tilde\omega)\mathbin{.}Tn\mathrm{d}t\\
&\qquad-aa\cos\varphi\mathbin{.}\cos(v-\tilde\omega)\mathbin{.}(\cos(v-\tilde\omega)+\cos E)Vn\mathrm{d}t\\
&\qquad+2aa\cos\varphi\mathbin{.}Vn\mathrm{d}t.
\end{aligned}
\]
\SourceLeaf{461}
\par Les deux dernières parties, qui contiennent le facteur \(Vn\mathrm{d}t\), deviennent, en mettant \(\dfrac{\cos(v-\tilde\omega)+e}{1+e\cos(v-\tilde\omega)}\) à la place de \(\cos E\):
\[
\begin{aligned}
&=\dfrac{aa\cos\varphi\mathbin{.}\sin(v-\tilde\omega)^2\mathbin{.}(2+e\cos(v-\tilde\omega))}{1+e\cos(v-\tilde\omega)}Vn\mathrm{d}t\\[4pt]
&=\dfrac{ar\sin(v-\tilde\omega)^2\mathbin{.}(2+e\cos(v-\tilde\omega))}{\cos\varphi}Vn\mathrm{d}t.
\end{aligned}
\]
\par Ainsi nous avons
\[
\begin{array}{rl}
68)\quad&\begin{aligned}[t]\mathrm{d}\tilde\omega&=(1-\cos i)\mathrm{d}\AscendingNode-aa\operatorname{cotang}\varphi\mathbin{.}\cos(v-\tilde\omega)\mathbin{.}Tn\mathrm{d}t\\
&\quad+\dfrac{ar\sin(v-\tilde\omega)\mathbin{.}(2+e\cos(v-\tilde\omega))}{e\cos\varphi}Vn\mathrm{d}t.\end{aligned}
\end{array}
\]
